% 本文件由 LaTeX 排版 Agent 根据有效 translation.json 自动生成。
% author=Tertullian; work_id=0309; title=灵魂的证言

\chapter{\WorkTitle{灵魂的证言}}

% structure: p0001 source=h1 target=omitted reason=page-title-already-rendered-by-parent

% structure: p0002 source=h2 target=section reason=relative-heading-rank; base=h2

\section{第一章}

如果有人为了从那些反对和迫害基督信仰真理之人的权威本身，来指控他们既自相矛盾又对我们不公，而决意从哲学家、诗人或这世界学识与智慧的其他大师的著作中，为这真理搜集证据，他就需要极其敏锐的精神和更大的记忆力来进行这项研究。事实上，我们的一些人，他们仍继续在古文学中进行探索性的劳作，并仍以此充满记忆，已经出版了我们手头现有的这类著作；在这些著作中，他们记述并证实了其传统之性质与来源，以及意见所依据的理由，从中可以立刻看出，我们所接受的并无新奇或怪异之处——没有什么是我们不能从普通和著名的著作中寻求支持的，无论是在从我们的信条中排除错误，还是将真理纳入其中。然而，人类心灵不信的硬心使他们轻视甚至自己的导师，尽管这些导师在其他方面受到认可并享有盛誉，只要他们触及用于辩护基督信仰的论点。于是，诗人们成了愚人，当他们用人的激情和故事来描述诸神时；于是，哲学家们成了没有理性的人，当他们叩击真理之门时。一个发表了几乎类似于基督徒情感言论的人，将在某种程度上被视为智者和有洞察力的人；而如果有人出于判断和智慧的一时做作，决意拒绝他们的仪式，或指控世界的罪恶，他必定会被打上基督徒的烙印。因此，我们将与那些在其最佳结果上被扭曲的文学和教导毫无关系，这些教导在错误中被相信，而非在其真理中。如果他们中的一些作者宣称有一位天主，且只有一位天主，我们也不会对此加以强调。不，就算承认在异教徒作家中没有任何基督徒认可的东西，以便使他无法说出任何一句责备的话。因为并非所有人都熟悉他们的教导；而那些熟悉的人，对其真实性也没有把握。人们更不会同意我们的著作，除非已经是基督徒，否则没有人会来寻求指导。我呼吁一个新的见证，是的，一个比所有文学都更知名、比所有教义都更被讨论、比所有出版物都更公开、比整个人更伟大的见证——我指的是所有属于人的东西。站出来吧，啊，灵魂，无论你是一个神圣而永恒的本质，正如大多数哲学家所相信的那样——如果是这样，你就越不可能说谎——还是你恰恰与神圣相反，因为确实是有死之物，正如伊壁鸠鲁独自认为的那样——那样的话，你在这种情况下就更没有诱惑去说假话：无论你是从天上领受的，还是从地里生出的；无论你是由数字组成的，还是由原子组成的；无论你的存在随着身体开始，还是你在后来才被放入其中；无论从何种来源，以何种方式，你使人成为一个理性存在，最高度地具有思想和知识的能力——站出来，给出你的见证。但我呼唤你，并非像在学校中塑造、在图书馆中训练、在雅典的学园和柱廊中餍足时那样，你吐露智慧。我呼唤你，简单的、粗野的、未受培养和未经教导的你，正如那些只拥有你的人们所拥有的你一样；完全是那路边的、街头的、作坊里的东西。我需要你的无经验，因为在你的小小经验中无人感到任何信赖。我向你要求那些你带入人身上的东西，那些你或是从你自己，或是从你的作者（无论他是谁）所知道的东西。你，正如我所深知，不是基督徒；因为人是成为基督徒的，而非生来就是基督徒。然而基督徒热切地要求你的见证；他们要求你，虽是一个外人，却要见证反对你的朋友们，好让他们在你面前感到羞愧，因为他们因那些证明你是同谋的事而仇恨和嘲笑我们。\par

% structure: p0004 source=h2 target=section reason=relative-heading-rank; base=h2

\section{第二章}

我们因宣告唯有一天主，唯祂堪当天主之名，万物由祂而来，祂是全宇宙的主，因而得罪人。若你确知这是真理，请作证；因为我们听见你公开而完全自由地（这是我们所未有的）在私下和公开场合呼喊：\Quote{愿天主赐予}和\Quote{若天主愿意}。藉此类言词，你宣告有一位独特的天主，并承认一切权能都属于祂，你作为臣民仰望祂的旨意。同时，你也否认其他神祇为真神，因你以它们各自的名字称呼它们——萨图尔努斯、朱庇特、马尔斯、弥涅耳瓦；因为你肯定唯祂是天主，你不给祂别的名字，只称天主；虽然你有时也称它们为神，但你显然是用一个并不真正属于它们的称谓，可以说是借来的称谓。我们所宣告的天主的性质，你并非不知：你常说\Quote{天主是善的，天主行善}；显然又暗示\Quote{但人是邪恶的}。你以反命题的方式，间接而形象地责备人的邪恶，因他离开了如此良善的天主。同样，如同在我们中间，作为慈爱良善的天主，\Quote{祝福}在我们的宗教和生活中是最神圣的行为，你也像基督徒所需的那样轻易地说\Quote{愿天主祝福你}；而当你把天主的祝福转为诅咒时，你的言语也同样与我们一同承认祂对我们拥有绝对而完全的权能。有些人虽不否认天主的存在，却认为祂既非监察者，也非治理者，也非审判者；他们特别鄙夷我们中因畏惧未来审判而归向基督的人，认为他们这样做是在尊重天主，使祂免于看守的操心与注意的烦劳——甚至不认为天主会发怒。因为他们说，若天主发怒，祂就会受败坏和激情的影响；但能受激情和败坏影响的事物也可能消亡，天主不能如此。然而这些人，在别处承认灵魂是神圣的、由天主赐予我们时，却撞上了灵魂本身的证言，这证言反驳了这些观点。因为灵魂若或是神圣的，或是天主所赐，无疑知道它的赐予者；若知道祂，无疑也畏惧祂，尤其因祂如此丰盛地赋予了它。它岂不畏惧祂吗？它如此渴望获得祂的恩宠，又如此焦虑避免祂的愤怒。那么，若天主不能发怒，灵魂对天主的自然畏惧从何而来？若无事物冒犯祂，如何能畏惧祂？所畏惧的岂非愤怒？愤怒从何而来，岂非从观察所行之事？什么导致警觉的看守，岂非将来的审判？审判从何而来，岂非来自权能？至高权威和权能属于谁，岂非唯独天主？所以，哦灵魂，你凭自己的知识，无人嘲笑你，无人阻止你，时刻准备呼喊：\Quote{天主洞悉一切}、\Quote{我把你交托给天主}、\Quote{愿天主报偿}、\Quote{天主必在我们之间审判}。这事如何发生，既然你不是基督徒？为何你头戴刻瑞斯的花冠，身披萨图尔努斯的紫袍，身着伊西斯女神的白衣，却呼求天主为审判者？站在埃斯科拉庇俄斯的雕像下，装饰朱诺的铜像，用暗色图形打扮弥涅耳瓦的头盔，你从未想到向这些神祇中任何一个呼求。在你的广场上，你向一位身在他处的天主呼求；你容许在你的神庙中向一位外来的神致敬。啊，真理的惊人证言，在魔鬼中间为我们基督徒获得了见证！\par

% structure: p0006 source=h2 target=section reason=relative-heading-rank; base=h2

\section{第三章}

但当我们说有魔鬼时——好像仅凭我们独自将魔鬼从人身上驱逐出去这一事实，并未证明它们的存在——某位克吕西波的弟子便开始撇嘴。然而你们的诅咒足以证明确有这些存在，并且你们深恶痛绝。为了表达对某人的强烈憎恨，你们称那以污秽、恶意、傲慢或任何我们归于恶灵的其他恶习困扰你的人为魔鬼。在表达恼怒、轻蔑或憎恶时，你们口中常挂着撒殚；这正是我们所认为的邪恶天使、谬误之源、全世界的败坏者，起初人因它而陷入违反天主诫命的境地。人因罪被交付于死亡，整个人类从他而出的后代都受到沾染，成为传递他定罪的渠道。你看，这便是你的毁灭者；虽然它只有基督徒或任何承认主的派别才完全知晓，但你们即使厌恶它，也已有几分认识它！\par

% structure: p0008 source=h2 target=section reason=relative-heading-rank; base=h2

\section{第四章}

即便如今，当此事关系到你的意见在某个更切合你的观点上，就它影响你个人的幸福而言，我们坚持认为：在生命逝去之后，你依然继续存在，并期待一个审判之日，并依照你的功过被判予痛苦或福乐，无论哪种方式都是永远的；并且，为能如此，你从前的实体必须回到你身上，即同一个人的物质和记忆：因为如果你不再被赋予那有感觉的身体组织，你就不能感受善或恶；若不呈现那承受审判痛苦的本人，就没有审判的依据。那基督宗教的观点，虽远比\OriginalTerm{毕达哥拉斯学派}{Pythagorean}的观点更为高贵，因为它不会将你转为禽兽；比\OriginalTerm{柏拉图学派}{Platonic}的观点更完整，因为它再次赐予你身体；比\OriginalTerm{伊壁鸠鲁学派}{Epicurean}的观点更值得荣耀，因为它使你免于消亡——然而，因着与它关联的名称，它却被视为不过是虚妄和愚蠢，且如所称，仅仅是一种僭妄。但是，若我们的僭妄被发现有你的支持，我们也不以此为耻。首先，当你谈论一位死者时，你会他说：\Quote{可怜的人}——可怜，当然不是因为他被夺去了生命的美善，而是因为他已被交付于惩罚和定罪。但在别的时候，你又提及死者摆脱了烦恼；你声称认为生命是一种负担，死亡是一种福分。你也常说到死者处于安息，当你带着食物和美味回到城外他们的坟墓时，你常是把祭品献给自己而非他们；或当你从坟墓回来时，你因酒力而蹒跚。但我想要你清醒的意见。当你表达自己的思想、远离他们时，你称死者为可怜。因为在他们宴席中，在某种意义上他们同在并与你一同躺卧，断不能指责他们的命运。你不得不奉承那些因他们的缘故你才如此丰盛宴乐的人。那么，你称那没有感觉的人为\Emph{可怜的}吗？你为何诅咒那个人，好像他能因你的诅咒而受苦，他的记忆带着某种旧伤的刺痛回到你身上？你的诅咒是：\Quote{愿大地重重压在他身上}，并且在死者之地有苦难\Quote{临到他的灰烬}。同样，在你对曾施恩于你的那个人所怀的善意中，你恳求\Quote{安息于他的骨骸和灰烬}，并期望他在死者中\Quote{享受愉快的安息}。如果你在死后没有受苦的能力，如果没有感觉残留——简言之，若与身体的分离就是你的虚无，什么使你对自己说谎，好像你还能在另一种状态受苦？不，你为何害怕死亡呢？若无物可感，死后就无可惧。因为虽可论说死亡之所以可怕，并非因为它以后威胁什么，而是因为它剥夺了生命的美善；但另一方面，由于它结束了生命的诸多不适——那些不适远为众多——死亡的恐惧被一种远超损失的收益所缓和。并且没有理由为失去美善之物而烦恼，这失去已被如此大的福分——从一切烦恼中获得解脱——充分补偿。那从一切可畏之物中解放出来的事物，并无可怕之处。若因生命经历甘甜而退缩，不愿放弃生命，那么，若不知道死亡的邪恶，就无需为死亡惊慌。你对它的恐惧证明你意识到它的邪恶。你绝不会认为它邪恶——你根本不会害怕它——除非你确信死后有使它成为邪恶、从而成为恐怖之事的东西。此刻我们暂时搁置那种害怕死亡的自然方式。任何人害怕不可避免之事，都是可怜之事。我转而从另一面，基于我们尘世生命期限之外的喜乐希望来论辩；因为对死后名声的渴望几乎是每个阶层与生俱来的。我没有时间谈论\Person{库尔提乌斯}和\Person{雷古卢斯}，或\Place{希腊}的勇敢的人们，他们向我们提供了无数为身后名声而轻视死亡的实例。如今谁没有这样的愿望：自己死后能常被记念？谁不竭尽全力，通过文学作品，或通过德行的单纯荣耀，甚至通过坟墓的辉煌，来保全自己的名声？灵魂的本性怎会怀有这些身后野心，并以如此惊人的努力预备那些只能在死后使用的事物？若未来对它全然未知，它就不会关心未来。但或许你以为，在脱离身体之后，你对自己继续感觉的肯定，超过对任何将来复活的肯定——这复活是被归咎于我们的教义，被视为我们僭妄的假设之一。但这教义也是灵魂本身的教义；因为若有人询问一个刚死之人，如同他仍活着，人立即会说：\Quote{他走了。}那么，他是被期待回来的。\par

% structure: p0010 source=h2 target=section reason=relative-heading-rank; base=h2

\section{第五章}

灵魂的这些见证是单纯的，正如真实的；是平常的，正如单纯的；是普遍的，正如平常的；是自然的，正如普遍的；是神圣的，正如自然的。我不认为它们对任何人会显得轻浮或无力，只要他思考大自然之威严，灵魂正是从大自然获得其权威。若你承认这位女主人的权威，也必然在门徒身上承认它。好，大自然在此是女主人，她的门徒是灵魂。但一方所教导或另一方所学习的一切，都来自天主——那位老师的老师。至于灵魂从其首席导师的教导可能知道什么，你可以从你内在的那一位来判断。思考那使你能够思考者；反思那在预感中是先知，在预兆中是占卜者，在将来事件中是预见者。若它（灵魂）作为天主赐予人的恩赐，懂得如何预言，岂非奇妙之事？若它认识那赐予它的天主，岂非极不寻常？即使它如此堕落，成为大敌阴谋的受害者，它仍不忘记它的创造主、祂的良善与法律，以及它自身及其仇敌的终极结局。那么，若它（灵魂）根源神圣，它的启示与天主赐给祂自己子民的知识相符，这有何独特呢？但那些不视灵魂的这些流露为先天本性的教导和内在知识的隐秘宝藏的人，会说这种说话的习惯及（可以说）恶习，是从广为流传的已出版书籍的观点中获得并得到确认的。毫无疑问，灵魂在文字之前存在，言语在书籍之前存在，思想在书写它们之前存在，人本身在诗人和哲学家之前存在。难道可以相信，在文学及其出版之前，没有我们指出过的那种话语出自人的口唇吗？难道没有人谈论天主和祂的良善，没有人谈论死亡，没有人谈论死者？我想，言语曾去乞讨；不，（主题仍然缺乏，没有它们言语甚至无法在今天存在，而今天言语已如此丰沛、丰富而有智慧；）如果那些如今如此容易提示、如此持续纠缠我们、如此贴近我们、仿佛生于我们唇上之事，在古代，在世界上任何文字存在之前——我以为，甚至在墨丘利存在之前——毫无存在，那么言语根本不可能存在。请问，文字本身从何得知并传播供言语使用那些从未有任何心智构思、任何口舌说出、任何耳朵听闻之事？但是，显然，既然天主的圣经，无论是属于基督徒或属于犹太人的（我们被接在其橄榄树上），远比任何世俗文学更为古老（或者，我们仅说，年代稍早，正如我们在适当之处证明其可信性时已表明）；若灵魂确实从著作中采取了这些话语，我们必须相信它是从我们的著作中采取，而非从你们的，因为它的教导更自然地来自较早的作品，而非较晚的。后者确实等候从前者获得教导；虽然我们承认光明曾从你们而来，但最初它仍是从第一源头流出的；我们将你们可能从我们这里取去并流传的一切，全数据为己有。既然如此，灵魂的知识是由天主还是由祂的书注入其中，便无关紧要。那么，人啊，你为何要坚持如此毫无根据的观点，认为灵魂的这些见证仅仅出自你们文学中的人性思辨，并且在日常使用中变得僵化呢？\par

% structure: p0012 source=h2 target=section reason=relative-heading-rank; base=h2

\section{第六章}

那么，相信你自己的著作吧；至于我们的圣经，更要相信神圣的著作；但在灵魂本身的见证上，给予大自然同样的信任。从这些中选择你所观察到的真理最忠实的朋友。若你自己的著作不被信任，则天主与大自然的见证不说谎。若你愿意相信天主与大自然，就相信灵魂；这样你便相信你自己。诚然，你重视灵魂，因为它给予你真正的伟大——你所归属的那一位；它是你的一切；没有它你既不能生也不能死；为了它你甚至将天主从你身边推开。既然你害怕成为基督徒，就把灵魂召到你面前，拷问她。她为何敬拜另一位？为何呼求天主的名字？为何当她意指该受诅咒的灵时，却提到魔鬼？为何她向着上天宣誓，而向大地发出日常的诅咒？为何她在一处服务，在另一处却呼求复仇者？为何她对死者作出判决？那些她拥有的基督徒短语是什么，尽管她既不愿见也不愿听基督徒？她为何将它们赐予我们，或从我们接受？她为何教导我们它们，或作为我们的学生学习它们？当实践如此不同时，要怀疑文字上的一致。否认普遍的本性——将这一点完全归因于我们的语言和希腊语（在我们看来如此相近），是十足的愚蠢。灵魂并非只对拉丁人和希腊人来自天堂的恩赐。人是属于地上万民的同一个名字：有一个灵魂和许多语言，一个精神与各种声音；每个国家有自己的言语，但言语的主题是共同的。天主无所不在，天主的良善也无所不在；魔鬼无所不在，对它们的诅咒也无所不在；呼求神圣审判也无所不在；死亡无所不在，对死亡的感知也无所不在；全世界都能找到灵魂的见证。没有一个人的灵魂，不从自身的光明中宣告那些我们不被允许高声说出的东西。因此，每个灵魂既是见证者，也是罪人：它为真理作证多少，错误的罪责就落在它身上多少；在审判之日，它将站在天主的法庭面前，无话可说。灵魂啊，你曾宣告天主，却不寻求认识祂；恶灵被你憎恶，却又是你崇拜的对象；地狱的惩罚被你预见，却未加避免；你曾有基督徒的气息，却同时是基督徒的迫害者。\par

\SourceRecord{Translated by S. Thelwall.；Ante-Nicene Fathers；Vol. 3.；Edited by Alexander Roberts, James Donaldson, and A. Cleveland Coxe.；Buffalo, NY: Christian Literature Publishing Co.,；1885.；<http://www.newadvent.org/fathers/0309.htm>.}
